\subsubsection{Mekanisme Pengiriman PGHD oleh Pasien}\label{subsubsec:pengiriman-pghd}


\begin{enumerate}
    \item \textbf{Pembangkitan kunci penandatanganan PGHD.}
    Pada saat registrasi akun pasien, selain pasangan kunci IOTA (\verb|iota_key_pair|) dan kunci proxy re-encryption (\verb|pre_public_key|) yang telah dibangkitkan oleh DecMed, klien pasien juga membangkitkan sepasang kunci publik dan kunci privat untuk penandatanganan PGHD menggunakan algoritma ECDSA, menghasilkan \verb|pghd_public_key| dan \verb|pghd_secret_key|. 
    ECDSA dipilih karena ukuran kuncinya yang lebih kecil namun menawarkan tingkat keamanan yang ekuivalen dengan RSA, sehingga lebih hemat komputasi untuk perangkat pasien.
    
    \item \textbf{Publikasi kunci publik ke IOTA.}
    Kunci publik \verb|pghd_public_key| dikirim dan disimpan pada ledger IOTA sebagai bagian dari metadata identitas pasien. 
    Dengan demikian, komponen lain di ekosistem DecMed dapat memverifikasi \textit{digital signature} PGHD menggunakan \verb|pghd_public_key| tanpa bergantung pada penyimpanan kunci di sisi klien.
    
    \item \textbf{Pembentukan payload plaintext dan \textit{inner signature}.}
    Ketika data PGHD akan dikirim, klien pasien terlebih dahulu membentuk payload PGHD dalam bentuk plaintext. 
    Klien kemudian menghitung nilai \textit{hash} plaintext dengan fungsi SHA-256,
    \[
        h_{\text{plain}} = \mathrm{SHA256}(\texttt{pghd\_plaintext}),
    \]
    dan menghasilkan \textit{digital signature} \verb|pghd_inner_signature| dengan menandatangani \verb|h_plain| menggunakan \verb|pghd_secret_key| melalui algoritma ECDSA\@. 
    Payload plaintext dan \verb|pghd_inner_signature| kemudian digabungkan ke dalam satu struktur data (field \verb|pghd_data| dan \verb|inner_signature|) sebelum dienkripsi.
    
    \item \textbf{Enkripsi payload PGHD.}
    Struktur data gabungan yang berisi payload PGHD dan \verb|pghd_inner_signature| kemudian dienkripsi menggunakan algoritma AES-GCM dengan kunci simetris acak \verb|aes_key| dan \verb|aes_nonce|, menghasilkan ciphertext \verb|enc_pghd|. 
    Selanjutnya, \verb|aes_key| dan \verb|aes_nonce| digabungkan dan dienkripsi menggunakan \verb|pre_public_key| milik pasien melalui mekanisme proxy re-encryption, menghasilkan pasangan \verb|enc_aes_key_nonce| dan \verb|capsule| serupa dengan mekanisme penambahan rekam medis oleh tenaga kesehatan.
    
    \item \textbf{Pembangkitan \textit{outer digital signature} atas ciphertext.}
    Untuk menjamin integritas dan \textit{non-repudiation} tanpa membebani komputasi \textit{on-chain}, klien pasien menghitung nilai \textit{hash} dari \verb|enc_pghd| (yang dapat berukuran besar) secara lokal di memori klien menggunakan fungsi SHA-256:
    \[
        h_{\text{cipher}} = \mathrm{SHA256}(\texttt{enc\_pghd}).
    \]
    Nilai \verb|h_cipher| kemudian ditandatangani menggunakan kunci privat pasien \verb|pghd_secret_key| dengan algoritma ECDSA, sehingga diperoleh \textit{outer digital signature} \verb|pghd_outer_signature|.
    
    \item \textbf{Pengunggahan off-chain dan pengiriman metadata ke \textit{smart contract}.}
    Klien pasien mengirimkan payload lengkap (\verb|enc_pghd|, \verb|h_cipher|, \verb|enc_aes_key_nonce|, \verb|capsule|, dan \verb|pghd_outer_signature|) ke proxy. 
    Proxy bertindak sebagai perantara \textit{off-chain} dengan mengunggah \verb|enc_pghd| ke kluster IPFS terlebih dahulu untuk mendapatkan \textit{Content Identifier} (CID). 
    Setelah CID diperoleh, proxy memanggil fungsi \texttt{submit\_pghd} pada \textit{smart contract} DecMed. 
    Payload yang dikirimkan ke jaringan IOTA\@: CID, \verb|h_cipher|, \verb|enc_aes_key_nonce|, \verb|capsule|, dan \verb|pghd_outer_signature|.
    
    \item \textbf{Verifikasi \textit{outer digital signature} dan penyimpanan state di IOTA.}
    Setelah menerima pemanggilan fungsi, \textit{smart contract} memproses metadata dengan langkah berikut:
    \begin{enumerate}
        \item memuat \verb|pghd_public_key| pasien dari state \textit{on-chain}, dan
        \item memverifikasi \verb|pghd_outer_signature| terhadap nilai \verb|h_cipher| yang diterima sebagai parameter menggunakan algoritma ECDSA\@.
    \end{enumerate}

    Jika verifikasi gagal, transaksi dibatalkan dan metadata tidak disimpan. Jika verifikasi berhasil, \textit{smart contract} secara permanen menyimpan kaitan antara CID IPFS, \verb|h_cipher|, dan metadata kriptografis lainnya pada ledger IOTA\@. 
    \textit{Smart contract} sama sekali tidak berinteraksi dengan data \verb|enc_pghd|.
\end{enumerate}